% !TEX program = lualatex
\documentclass[aspectratio=43,professionalfonts,handout]{beamer}
\usepackage{beamer_template}

\newcommand{\LectureNum}{5}
\newcommand{\LectureTitle}{対数（１）}
\newcommand{\TextPages}{p.\,112-114}
\newcommand{\NextTopic}{対数（２）・底の変換}
\newcommand{\NextPages}{p.\,114-115}
\newcommand{\CourseName}{基礎数学B}
\renewcommand{\TermName}{後期}

\begin{document}

% -----------------------------------------------
% スライド1: 表紙＋目標＋用語・定義
% -----------------------------------------------
\begin{frame}{\CourseName\quad 第\LectureNum 回「\LectureTitle 」}
  {\small \TermName\quad \TeacherName\quad ／\quad 教科書 \TextPages\quad
  \textbf{目標}：対数の定義と基本性質を理解し計算できる}

  \begin{block}{用語・定義}
  \begin{description}[labelwidth=5.5em]
    \item[\kw{対数の定義：$\log_a b = c \Leftrightarrow a^c = b$}] $a > 0$, $a \neq 1$, $b > 0$ として $a^c = b \Leftrightarrow c = \log_a b$
    \item[\kw{$\log_a 1 = 0$、$\log_a a = 1$}] 底 $a$ を0乗すると1，1乗すると $a$ になることから
    \item[\kw{真数条件：$b > 0$、底の条件：$a > 0$, $a \neq 1$}] 対数 $\log_a b$ は $b > 0$ のときのみ定義される
  \end{description}
  \end{block}
\end{frame}

% -----------------------------------------------
% スライド2: 公式・性質
% -----------------------------------------------
\begin{frame}{公式・性質}
  \begin{block}{}
  \begin{itemize}
    \item $\log_a MN = \log_a M + \log_a N$
    \item $\log_a (M/N) = \log_a M - \log_a N$
    \item $\log_a M^k = k \cdot \log_a M$
  \end{itemize}
  \end{block}
\end{frame}

% -----------------------------------------------
% スライド3: 例1→問1
% -----------------------------------------------
\begin{frame}{例1→問1}
  \begin{exampleblock}{例1) p.113（対数の値を求める）}
    （板書で解説）
  \end{exampleblock}

  \vfill

  \begin{block}{問1) p.113\quad 自分で解いてみよう}
    （ノートに解くこと）
  \end{block}
\end{frame}

% -----------------------------------------------
% スライド4: 例2→問2
% -----------------------------------------------
\begin{frame}{例2→問2}
  \begin{exampleblock}{例2) p.114（対数の性質を使った計算）}
    （板書で解説）
  \end{exampleblock}

  \vfill

  \begin{block}{問2) p.114\quad 自分で解いてみよう}
    （ノートに解くこと）
  \end{block}
\end{frame}

% -----------------------------------------------
% まとめ
% -----------------------------------------------
\begin{frame}{まとめ}
  \begin{itemize}
    \item 対数は指数の「逆」：$a^c = b \Leftrightarrow \log_a b = c$
    \item $\log_a b + \log_a c = \log_a(bc)$，$\log_a b - \log_a c = \log_a(b/c)$
    \item 計算前に必ず真数条件 $b > 0$ と底の条件を確認する
  \end{itemize}

  \vfill
  {\small 基本問題：225, 226, 227, 228}
\end{frame}

\end{document}
